% ------------------------------------------------------------------------
% file `travail_de_vacances-exercise-10-solution-body.tex'
%   in folder `xsim/'
%
%     solution of type `exercise' with id `10'
%
% generated by the `solution' environment of the
%   `xsim' package v0.20c (2021/02/03)
% from source `travail_de_vacances' on 2023/07/04 on line 568
% ------------------------------------------------------------------------
    \begin{tasks}
        \task La première quinzaine elle a consommé $\dfrac{3}{4}$ de son
        abonemment, il lui reste donc $\dfrac{1}{4}$. La deuxième quinzaine,
        elle a consommé $\dfrac{2}{5}$ de ce qu'il lui restait, c'est-à-dire,
        \begin{equation*}
            \dfrac{2}{5}\cdot\dfrac{1}{4} = \dfrac{1}{10}
        \end{equation*}
        Elle aura donc consommée durant tout le mois,
        \begin{equation*}
            \dfrac{3}{4} + \dfrac{1}{10} = \dfrac{17}{20}
        \end{equation*}

        \task Il lui reste $\dfrac{3}{20}$ non consommé.
        \task En utilisant la règle de trois,
        \begin{align*}
            \dfrac{3}{20}  & \longrightarrow \SI{12}{min} \\
            \dfrac{1}{20}  & \longrightarrow \SI{4}{min}  \\
            \dfrac{20}{20} & \longrightarrow \SI{80}{min}
        \end{align*}
        Au début du mois de juin \SI{80}{min} sont disponibles.
    \end{tasks}
